\section{Practical Design Guidance and Broader Context}\label{app:practical_guidance}

\subsection*{Setting the ROPE}

Three design inputs must be fixed before data collection: the ROPE, the precision goal
$\omega_{\rm goal}$, and the maximum sample size $N_{\rm max}$.
The precision goal should satisfy $\omega_{\rm goal} \leq \Delta_{\rm ROPE}$ and directly
determines the expected stop iteration $N_{\rm goal}$
(Equation~\ref{eq:pitg_stop_iteration}); $N_{\rm max}$ is set by budget.
Of these three subjective inputs the ROPE is the most consequential.
It should be determined from domain knowledge rather than from the data,
a requirement that forces the researcher to commit to a practically meaningful
effect-size threshold before observing outcomes \citep{kruschke2018}.
In clinical research the Minimal Clinically Important Difference (MCID) provides a
principled basis \citep{jaeschke1989}; in bioequivalence testing, regulatory guidelines define the equivalence
margin \citep{schuirmann1987}; in industrial quality
control, engineering tolerances provide an equivalent principled basis.
When no domain standard exists, sensitivity analyses across a range of ROPE widths
are recommended before committing to a single value \citep{lakens2018}.

\subsection*{Broader Adoption Considerations}

Three considerations support broader adoption of the DPitG framework.

First, formally accepting the null via the HDI+ROPE framework's criterion parallels the objective
of frequentist equivalence testing through the Two One-Sided Tests (TOST) procedure
\citep{schuirmann1987, lakens2018}, providing a Bayesian counterpart to a methodology
established in regulatory and applied settings \citep{fda2010bayesian}.

Second, because $\omega_{\rm goal}$, the ROPE, and $N_{\rm max}$ must all be specified
before data collection begins, DPitG is structurally compatible with pre-registration
standards \citep{nosek2018}, directly reducing the researcher degrees of freedom,
the analytic choices that enable post-hoc optimisation of results, that drive
selective reporting.

Third, the replication crisis has been partly attributed to underpowered studies whose
imprecise estimates fail to replicate \citep{osc2015, gelman2014}; DPitG addresses this
at the design stage by making precision a non-negotiable stopping criterion rather than
an aspiration, one that is structurally enforced by construction.
